% Exact scope of the relative-product argument for task A.7.
% Attribution: AI-produced by GPT-6 Astra (OpenAI, ChatGPT Pro mode).

\subsection{What relative products do and do not prove about the pentagon}

\paragraph{LX.1 (Proved): relative products preserve a specified pentagon equation.}
Let $P_4=\widehat\pi_1(M_{0,4})=F$ and
$P_5=\widehat\pi_1(M_{0,5})$. Fix the five continuous
substitution maps $e_j:F\to P_5$ occurring in the group
pentagon, with all paths and product order fixed. Let
$N_5\triangleleft P_5$ be open, and let $N_4\triangleleft F$
be an open subgroup satisfying $e_j(N_4)\subseteq N_5$ for
all five maps. Such $N_4$ exists: intersect the five inverse
images and take a characteristic open subgroup contained in
that intersection. Write
\[
 D_i=P_i/N_i,\quad P_{i,p}=P_i/O^p(N_i),\quad
 P_{i,\rm nil}=P_i/R_{\rm nil}(N_i).
\]
Every $e_j$ induces continuous maps in these three categories,
and the actual quotient isomorphisms
\[
 P_{i,\rm nil}\simeq\prod_{p,D_i}P_{i,p}
\]
intertwine them. Consequently a parameter in the relative
pronilpotent source satisfies the pentagon there if and only if
all its relative pro-$p$ components satisfy the corresponding
pentagon, with their common finite-base coordinates retained.

\emph{Proof.}
For any homomorphism $e:N_4\to N_5$, composition with a finite
$p$-group quotient of $N_5$ proves
$e(O^p(N_4))\subseteq O^p(N_5)$; the same argument for finite
nilpotent quotients gives the analogous residual inclusion.
This proves the induced maps. The product comparison for any
profinite group with an open normal subgroup follows from
\[
 N/R_{\rm nil}(N)\simeq\prod_pN^{(p)}:
\]
finite nilpotent groups are the products of their Sylow groups,
and coprime Goursat together with compactness gives surjectivity.
For the extension over $D_i$, lift the common $D_i$ coordinate
once, then correct its tuple by an element of the displayed
kernel product. This proves the stated isomorphism without a
chosen splitting.

Word evaluation commutes with every continuous homomorphism.
If $\mathcal P(f)$ denotes the ordered pentagon defect, its
projection to each component is therefore exactly
$\mathcal P_p(f_p)$. The product is injective, so
$\mathcal P(f)=1$ in $P_{5,\rm nil}$ if and only if all these
components are one. The same observation applies to additional
two-cycle and three-cycle equations, provided their substitution
maps are included when choosing compatible finite bases.
Automorphisms of a whole quotient diagram also glue
componentwise when they induce one common finite-base diagram
automorphism. This is a statement about the diagram quotient;
it does not assert that arbitrary lists of four-point
automorphisms already define a five-point diagram automorphism.
\hfill$\square$

\paragraph{LX.2 (Exact unproved step): lifting the restricted solution space.}
The identity in LX.1 does not supply the epimorphism
\[
\begin{gathered}
 \{(1,f)\in\widehat{\mathrm{GT}}:\text{specified retained data}\}\\
 \longrightarrow
 \{\text{compatible relative pentagon solutions}\}.
\end{gathered}
\]
Neither R1 nor R2 states this epimorphism. Their proven
finite transport theorem retains a triple of peripheral
conjugators and the one peripheral product relation in a free
pants or modular source. Adding the pentagon requires a new
finite transport theorem for the five correlated substitutions
of the same parameter $f$, or an equivalent theorem for
automorphisms of the compatible four- and five-point diagram.
No such theorem follows from the supplied hypotheses alone.

\emph{Explanation with the exact logical boundary.}
The count of D.50.11 treats independent boundary values in
specified classes and extracts the coefficient of their product
being one. Its surface extension D.50.12 does the same for
\[
 \prod_{i=1}^g[a_i,b_i]\,c_1\cdots c_b=1.
\]
The pentagon instead requires an ordered relation between
$e_1(f),\ldots,e_5(f)$, all evaluations of one unknown $f$.
These five values are correlated by the substitution maps.
The independent-boundary character sum cannot silently be
restricted to that correlated subset; the Schur scalarization
and subsequent Mobius inversion count different sets.

There is also a distinct source issue. A five-punctured sphere
has free fundamental group of rank four. The group
$\pi_1(M_{0,5})$ appearing in the pentagon is the fundamental
group of a configuration space, with the split forgetting
sequence $F_3\rtimes F_2$ and abelianization of rank five.
Thus applying the genus-zero, five-boundary formula is not
applying a presentation of $P_5$. One cannot replace the
configuration group by the fiber curve merely because both
involve five marked points.

Even lifting the two free groups separately in the forgetting
sequence leaves the essential compatibility condition. If
$\rho:F_2\to\operatorname{Aut}(F_3)$ is its monodromy action,
candidate lifts $\beta,\gamma$ must satisfy
\[
 \beta\rho(b)\beta^{-1}
   =\operatorname{Inn}(h_b)\rho(\gamma(b))
       \quad(b\in F_2),
\]
with the appropriate continuous cocycle law for $h_b$.
R1 retains peripheral conjugators, but imposes no such
monodromy compatibility. LU.4 provides an independent example
of precisely why an extra commutation requirement can obstruct
a lift even when a quotient automorphism is perfectly valid.

LP.5 gives actual full gentle elements failing the pentagon,
so satisfying all the peripheral and generation conditions
used by R1 does not imply that additional relation. This is
an explicit obstruction to that inference, not a counterexample
to the still stronger surjectivity assertion displayed above.
In particular nothing here proves that the R2 mixed-sign
element itself fails the pentagon. Its absolute pro-prime
images are identity, so the degree-seven and degree-nine
detectors do not decide that element.

\paragraph{LX.3 (Proved): the exact compactness criterion.}
Let $\mathcal C=\{\Phi\in\widehat{\mathrm{GT}}_0:
\Delta(\Phi)=1\}$, and let
$\mathcal P:\widehat{\mathrm{GT}}_0\to P_5$ be the continuous
pentagon-defect map in the fixed convention. For each open
normal $J\triangleleft P_5$, set
\[
 \mathcal C_J=
 \{\Phi\in\mathcal C:\mathcal P(\Phi)\in J\}.
\]
Then
\[
 \widehat{\mathrm{GT}}\cap\mathcal C\ne\varnothing
 \quad\Longleftrightarrow\quad
 \mathcal C_J\ne\varnothing\ \hbox{for every open normal }J.
\]
Equivalently, vanishing of $\Delta$ on the full group is
witnessed by the emptiness of one exact finite pentagon test
on actual gentle residues. Arbitrary finite candidates that
have not been shown to lift to the gentle group are not such
residues.

\emph{Proof.}
The sets $\mathcal C_J$ are closed in the compact group
$\widehat{\mathrm{GT}}_0$. They are directed under finite
intersection because
$\mathcal C_{J_1\cap J_2}=\mathcal C_{J_1}\cap\mathcal C_{J_2}$.
If each is nonempty, compactness supplies an element of their
intersection. Its defect lies in every open normal subgroup
of $P_5$, hence is one. The reverse direction is immediate.
For a fixed $J$, the map $\mathcal P\bmod J$ and $\Delta$
are continuous finite-valued maps, so both factor through
one finite quotient of the compact profinite parameter group.
This proves the finite formulation. If one instead uses raw
finite pairs, the compactness argument must also impose all
finite automorphism and gentle relations simultaneously and
prove nonemptiness after every finite common refinement.
\hfill$\square$

\paragraph{LX.4 (Scope of the present result).}
The underlying relative group product holds at every fixed
compatible finite base, including mixed two- and three-primary
bases; LX.1 proves the coordinatewise assertion with the
pentagon. A smallest five-point level detecting $\Delta$ has
not been identified here, and no surviving shadow proves a
full mixed-sign element. The unresolved mathematical statement
is the restricted lifting theorem of LX.2, or an independent
finite obstruction as in LX.3. Calling the product decomposition
itself false would overstate what has been established.
